\chapter{Dataclass e modelagem de dados}

\section{Por que modelar antes de inserir no banco?}

Ao parsear XML, voce recebe atributos como strings soltas. Se voce passar dicionarios por todo o projeto, o codigo fica fragil:

\begin{lstlisting}[style=devbutterpython]
record["startDate"]
record["value"]
record["sourceName"]
\end{lstlisting}

O problema e que erros de chave aparecem tarde. Alem disso, nao fica claro quais campos existem.

Com \code{dataclass}, voce cria uma estrutura explicita:

\begin{lstlisting}[style=devbutterpython]
record.start_date
record.value_numeric
record.source_name
\end{lstlisting}

\section{Modelo HealthRecord}

Arquivo sugerido: \code{src/models/health_record.py}

\begin{lstlisting}[style=devbutterpython]
from dataclasses import dataclass


@dataclass(slots=True)
class HealthRecord:
    apple_record_id: int
    type: str | None
    source_name: str | None
    source_version: str | None
    device: str | None
    unit: str | None
    creation_date: str | None
    start_date: str | None
    end_date: str | None
    value: str | None
    value_numeric: float | None
    metadata: dict[str, str]
\end{lstlisting}

\section{Explicando cada nomezinho}

\begin{itemize}
  \item \code{@dataclass}: decorator que gera automaticamente metodos como \code{__init__} e \code{__repr__}.
  \item \code{slots=True}: reduz overhead de memoria removendo o \code{__dict__} dinamico de cada instancia.
  \item \code{apple_record_id}: id sequencial criado pelo seu parser, diferente do id do banco.
  \item \code{type}: tipo do Record do Apple Health, por exemplo \code{HKQuantityTypeIdentifierHeartRate}.
  \item \code{source_name}: origem do dado, por exemplo Apple Watch ou iPhone.
  \item \code{value}: valor original em texto.
  \item \code{value_numeric}: tentativa de conversao para numero.
  \item \code{metadata}: informacoes extras de \code{MetadataEntry}.
\end{itemize}

\section{Por que manter value e value\_numeric?}

Nem todo valor do Apple Health e numerico. Alguns sao categorias.

Numericos:

\begin{verbatim}
92
1200
75.4
\end{verbatim}

Categoricos:

\begin{verbatim}
HKCategoryValueSleepAnalysisAsleepCore
HKCategoryValueSleepAnalysisInBed
\end{verbatim}

Por isso salvamos os dois:

\begin{itemize}
  \item \code{value}: preserva o valor original;
  \item \code{value_numeric}: permite agregacoes numericas quando possivel.
\end{itemize}

\section{Funcao safe\_float}

\begin{lstlisting}[style=devbutterpython]
def safe_float(value: str | None) -> float | None:
    if value is None:
        return None

    try:
        return float(value)
    except ValueError:
        return None
\end{lstlisting}

Ela tenta converter texto para numero. Se nao conseguir, retorna \code{None}. Isso evita quebrar o parser quando encontrar valores categoricos.

\section{Exercicio}

Implemente testes:

\begin{lstlisting}[style=devbutterpython]
def test_safe_float_with_number():
    assert safe_float("92.5") == 92.5


def test_safe_float_with_none():
    assert safe_float(None) is None


def test_safe_float_with_text():
    assert safe_float("HKCategoryValueSleepAnalysisAsleepCore") is None
\end{lstlisting}
